\paragraph{The component Bergman bound and its volume normalization.}
Fix one good complex embedding of the coefficient field. Write
\(f=[E:\mathbb Q]/2\), \(\Delta=|\operatorname{disc}E|\),
\(R_E=\Delta^{1/(2f)}\), \(r\ge2\) for the Hermitian rank,
\(d\) for the positive host integer and \(M\) for the deep level.
Let \(w\ge1\) bound the sum, over the CM type, of the spreads
\(\max_i\lambda_{\tau i}-\min_i\lambda_{\tau i}\) of the flag weights.
The complex dimension of the product of full flag domains is
\(b=fr(r-1)/2\). The constants in the component estimate are
\begin{align*}
 \rho&=[200r^3(w+f)]^{-1},&
 \varepsilon&=\min\{1,\pi/[2M\sqrt d\,(fw+2)\log2]\},\\
 A&=(1+2(2fr)^2\Delta/\varepsilon)^{2fr^2},&
 V_*&=(15M)^{2fr^2}[3(2f)^2r^4R_E^2]^{6fr^2},\\
 C_{\rm loc}&=\frac{8\,b!}{(\pi\rho^2)^b}\,A,&
 \beta_*&=C_{\rm loc}V_* .
\end{align*}
Here \(V_*\) is the upper bound called \(\mathrm{massLev2}\) in
the Tamagawa comparison. The normalizing mass itself is the sum
of the volumes of the chosen arithmetic fundamental sets.
\src{src/Development/AbelianVarietyTheory/Definitions/Divisors/UnifOfSpaceEq.lean}{103}{Flag radius, return bound and local factor};
\src{src/Development/WeilBound/RepresentationTheoryGL/TubeMeasures/DisjointTubes/MassLev2.lean}{13}{Level-dependent total mass bound};
\src{src/Development/WeilBound/RepresentationTheoryGL/TubeMeasures/DisjointTubes/ConstantSupremum/SupremumConstantLogBound.lean}{14}{Product defining the normalized supremum constant}.

\begin{lemma}[Evaluation on a flag ball]\label{bergvol:local}
At an orthonormal frame \(x\), let \(B_x\) be the lower unipotent
flag chart of Euclidean radius \(\rho\). Write \(d\mu\) for the flag
measure specified below. For every holomorphic flag section \(P\),
\[
        |P(x)|_{\rm Pet}^2\le
        \frac{2b!}{(\pi\rho^2)^b}\int_{B_x}|P|_{\rm Pet}^2\,d\mu .
\]
\end{lemma}
\begin{proof}
In the frame \(x\) each factor has a diagonal Hermitian matrix
\(J\), with the positive signs first and the negative signs last.
Write the chart matrix as \(U=I+Z\) with \(Z\) strictly lower
triangular; the entries over all factors are the Euclidean coordinates.
In particular \(\|Z\|_{\rm op}\le\|Z\|_{\rm HS}\le\rho\).
Let \(D_{\tau i}\) be the absolute determinant of the leading
\(i\)-by-\(i\) Gram block of \(U^*JU\). The endpoint determinants
\(D_{\tau0},D_{\tau r}\) equal one. On the chart the squared norm
and invariant measure in question are
\[
 |P|_{\rm Pet}^2=|P(xU)|^2
 \prod_{\tau}\prod_{i=0}^{r-1} (D_{\tau,i+1}/D_{\tau i})^{-\lambda_{\tau i}},
 \qquad
 d\mu=\left(\prod_{\tau}\prod_{i=0}^{r-1}D_{\tau i}^{-2}\right)dV .
\]
The difference of each Gram block from the corresponding sign matrix
has operator norm at most \(2\rho+\rho^2\le3\rho\). Its singular
values therefore lie in \([1-3\rho,1+3\rho]\), and the block has
the required inertia by continuity. Thus the whole ball belongs to
the flag domain and \(|\log D_{\tau i}|\le6r\rho\).
Subtract the midpoint of the extreme weights from each row of
weights at \(\tau\). This leaves the logarithm of the norm factor
unchanged by telescoping the endpoint determinants. Its absolute value
is at most \(6r^2\rho w\le1/20\); the absolute logarithm of the
density is at most \(12fr^2\rho\le1/10\). At the origin both
factors equal one. The holomorphic mean-value inequality on the
Euclidean ball gives the assertion since its volume is
\((\pi\rho^2)^b/b!\) and \(\exp(1/20+1/10)<2\).
\end{proof}

\begin{lemma}[Overlaps at a frame with no short lattice vector]\label{bergvol:returns}
Suppose the arithmetic group preserves a Hermitian \(O_E\)-lattice
\(\Lambda'\subset E^r\) self-dual for the \(O_E\)-valued form.
Let \(\|v\|_x^2\) be the sum of the positive majorants of that form
over the fixed CM type, obtained by changing the negative signs to
positive in \(x\). If its minimum on \(\Lambda'\setminus\{0\}\)
is at least \(\varepsilon\), at most \(A\) arithmetic elements
return \(B_x\) to itself, meaning \(\gamma B_x\cap B_x\ne\varnothing\).
\end{lemma}
\begin{proof}
Set \(m=2fr\), the real lattice dimension. The Euclidean dual of the
embedded lattice is \(2J\) times its dual for
\(\operatorname{Tr}_{E/\mathbb Q}h\), since that trace is twice
the real inner product after applying the sign involution \(J\).
The trace dual lies in \(\Delta^{-1}\Lambda'\): testing against
\(O_E\Lambda'\) puts each Hermitian pairing in the codifferent,
which is carried into \(O_E\) by \(\Delta\), as seen by inverting
the integral trace matrix. Hence the Euclidean dual minimum in length
is at least \(2\sqrt\varepsilon/\Delta\).
The Euclidean transference inequality
\(\lambda_m(L)\lambda_1(L^*)\le m\) applies to this full rank lattice.
It gives \(m\) independent vectors of length at most
\(m\Delta/(2\sqrt\varepsilon)\). They span \(E^r\) over \(E\):
otherwise their real span would have dimension less than \(m\).
Select an \(E\)-basis from them.
Banaszczyk,
\emph{New bounds in some transference theorems in the geometry of numbers},
\S2~\cite{banaszczyk}.

In the chart the positive subspace is the graph of a matrix \(T\)
with \(\|T\|\le\rho/(1-\rho)\), by inverting its upper block.
Its majorant is between one half and twice the central squared norm.
Indeed singular value decomposition reduces its change to independent
positive-negative planes with graph slope \(t\le\|T\|\); the
eigenvalues of the new majorant there are
\((1+t)/(1-t)\) and its inverse, by using the orthogonal lines
spanned by \((1,t)\) and \((t,1)\). Unpaired directions retain
eigenvalue one. A returning \(\gamma\) carries one chart majorant
isometrically onto another; hence \(\|\gamma v\|_x\le2\|v\|_x\).
All images of the selected basis vectors have length at most
\(m^2\Delta/\sqrt\varepsilon\), allowing the extra factor \(m\ge1\).
Lattice points in a ball of radius \(T_0\) number at most
\((1+2T_0/\sqrt\varepsilon)^m\), by packing disjoint open balls of
radius \(\sqrt\varepsilon/2\). The \(r\) basis images determine
\(\gamma\). Taking the \(r\)-th power of this count yields \(A\).
\end{proof}

\begin{proposition}[From cusp and mass data to a normalized kernel]\label{bergvol:global}
Assume the following on every analytic flag component at the chosen
embedding.
\begin{enumerate}
\item The automorphic sections under consideration are bounded in
Petersson norm. For each nonzero such cusp section there is an
orthonormal frame \(x\) where its norm is at least half the component
supremum and the minimum in Lemma~\ref{bergvol:returns} is at least
\(\varepsilon\), for the same \(\Gamma\)-stable self-dual lattice.
\item There are invariant flag measures in the chart normalization
of Lemma~\ref{bergvol:local} and measurable arithmetic fundamental
sets \(F_j\). With \(V=\sum_j\mu_j(F_j)\), assume \(0<V\le V_*\).
The pushforward of the adelic probability measure on arithmetic
and frame-invariant functions is \(\sum_j d\mu_j|_{F_j}/V\).
\end{enumerate}
Then for each cusp section under consideration the global squared supremum
is bounded by \(\beta_*\) times its squared Petersson norm for that probability. On any nonempty
algebraic component \(Z\) which corresponds to a union of these
analytic components of positive mass, the same bound holds with
the conditional probability on \(Z\). In particular the scalar
evaluation Bergman kernel on \(Z\) is at most \(\beta_*\).
\end{proposition}
\begin{proof}
At the chosen half-maximum frame, project the flag ball onto the
arithmetic quotient. Each fiber in the ball has at most \(A\)
points, since mapping one of them to the others requires elements
which return the ball. Equivalently, cut a measurable fundamental
domain into its arithmetic translates and integrate its invariant
integrand there; the number of copies of a point in the ball is
bounded by the return count, including stabilizers if needed.
Lemma~\ref{bergvol:local} and the half-maximum then give
\(\sup_{x\in F_j}|P(x)|_{\rm Pet}^2\le C_{\rm loc}
\int_{F_j}|P|_{\rm Pet}^2\,d\mu_j\).
Adding the nonnegative integrals and inserting \(V\) proves the
global result. For \(Z\) use the identical calculation summing only
over its analytic components and multiplying the conditional
integral by \(V_Z=\sum_{j\subset Z}\mu_j(F_j)\le V_*\).
Taking the supremum of the evaluation ratio over sections gives
the Bergman bound. Its integral against the conditional probability
bounds the rank as well: in any finite orthonormal system the sum
of the squared pointwise norms equals the squared norm of scalar
evaluation on its span, and each summand integrates to one.
\end{proof}

For clarity, the half-maximum hypothesis admits a quantitative test
which isolates the cusp input. Put \(H=2M\sqrt d\). If at a
half-maximum frame a lattice vector satisfies \(\|v\|_x^2<
\varepsilon\le1\), every conjugate of the algebraic integer
\(h(v,v)\) has absolute value less than one, since it is bounded
by the corresponding majorant. Thus \(h(v,v)=0\) by its norm.
For each such isotropic vector use the flow parameter scaled by
\(k_0=2|\tau(\xi_0)|\), independent of \(\tau\).
In this parameter the input needed is a transvection
period \(0<H'\le H\), acting with trivial scalar multiplier, and its
holomorphic one-variable cusp restriction \(G(\zeta)\). In the
normalization with scalar metric factor one at the center it
suffices that \(G(0)=P(x)\), that \(G\) be \(H'\)-periodic and
holomorphic on a neighborhood of \(\Im\zeta\ge-t\) for
\(t=1/(2\varepsilon)\), and that it tend uniformly to zero as
\(\Im\zeta\to+\infty\). The flag metric on the flow gives
\[
       |G(\zeta)|\le 2\cdot 2^{w/2}|G(0)|
                     \qquad(\Im\zeta=-t).
\]
To verify the metric assertion, write \(u\) for the coordinates of
\(\tau v\) at the orthonormal central frame, with positive columns
first. Since \(\tau\kappa_0=-i k_0\) on the chosen CM type, the
flow in the scaled parameter is \(U=I-i\zeta uu^*J\).
Isotropy gives \(u^*Ju=0\) and
\(U^*JU=J+2\Im\zeta\,Juu^*J\).
For \(S_i=\sum_{a=1}^i J_{aa}|u_a|^2\), the rank-one determinant
formula on each leading block now gives
\(D_{\tau i}=|1+2\Im\zeta\, S_i|\).
The \(S_i\) increase from zero to \(\|u\|^2/2\), then decrease to
zero, since the total positive and negative sums coincide.
Here \(\|u\|^2\) is the local majorant, less than
\(\varepsilon\) by hypothesis. All signed determinants retain their
central signs for \(\Im\zeta>-1/\varepsilon\). Thus the flow stays
inside the flag domain in this entire half-plane and the restriction
of a holomorphic section is holomorphic there.
At \(\Im\zeta=-1/(2\varepsilon)\), the \(D_{\tau i}\) decrease from
one to a number greater than \(1/2\), then increase to one.
The total variation of \(\log D_{\tau i}\) is at most \(2\log2\).
Centering the weights at their midrange in the squared metric formula
of Lemma~\ref{bergvol:local} bounds its absolute logarithm by
\(\sum_\tau\operatorname{spread}(\lambda_\tau)\log2\le w\log2\).
The boundary raw coefficient is therefore bounded by \(2^{w/2}\)
times its Petersson norm, which itself is at most \(2|G(0)|\).

For \(y=\Im\zeta\ge0\), the same variation argument gives a raw
coefficient of size at most \(2|G(0)|(1+\varepsilon y)^{w/2}\).
Thus for \(q=\exp(2\pi i\zeta/H')\) the growth is at most a polynomial
in \(\log(1/|q|)\). By Cauchy's estimate on circles shrinking to
zero, every negative Laurent coefficient vanishes. It suffices to
show that the constant Laurent coefficient vanishes, equivalently
that the average of this restriction along one real period is zero
at any fixed height.

For the analytic cusp space of this construction the constant is
indeed zero. Its definition imposes bounded Petersson norm and,
for every nonzero isotropic \(\xi\in\Lambda\) and every frame,
zero average of the coefficient on the following \(f\)-dimensional
real torus. Put \(F=E^+\); its period lattice in the raw flow
parameters is \(L=m e_\chi\,\iota(O_F)\), where \(\iota\) is the
real embedding vector, \(m\) is the integer transvection period of
the flag datum and \(e_\chi\) kills its finite order multiplier.
This is a full lattice by the nonzero discriminant of \(F\).
The real period in direction \(\iota(a)\) follows from the
rational transvection for \(a\in O_F\) and the multiplier is
trivial after its \(e_\chi\)-th power. Both boundedness and
zero average persist under complex linear combinations, so they
hold on the full span used for the analytic Bergman kernel.
\src{src/ReciprocityTheory/ReductiveGroups/ArithmeticGeometry/FundamentalDomains/ModerateCoefficientPositivity.lean}{58}{Definition of the analytic cusp space and its zero Fourier coefficient}.

Choose an integer \(c_1>0\) with \(\xi=c_1v\in\Lambda\), possible
since \(\Lambda\) is a full lattice in the rational vector space
underlying \(E^r\). Apply the flow in independent raw
parameters \(z_\tau\) to \(\xi\), and call the resulting holomorphic
coefficient \(A(z)\). For \(y_\tau=\Im z_\tau\ge0\), the same
principal minors are \(1+2 k_0 y_\tau S_i\), now with \(S_i\)
computed for \(\tau\xi\). Their variation and bounded Petersson norm
give
\[
 |A(t+iy)|\le C_P\prod_\tau
 (1+k_0 y_\tau\|\tau\xi\|_x^2)^{\operatorname{spread}(\lambda_\tau)/2}
 \qquad(t\in\mathbb R^f).
\]
Here the norm is the single-place majorant. In particular growth is
at most polynomial when one imaginary coordinate tends to infinity
and the others stay zero. For \(\eta\in L^*\) let \(a_\eta(y)\)
be the Fourier coefficient of \(t\mapsto A(t+iy)\) using
\(\exp(-2\pi i\eta\cdot t)\) and probability Haar measure on the
real torus. Holomorphy and integration by parts on this torus give
\(\partial_{y_j}a_\eta=-2\pi\eta_j a_\eta\), hence
\(a_\eta(y)=e^{-2\pi\eta\cdot y}a_\eta(0)\).
The polynomial bound implies \(a_\eta(0)=0\) if any \(\eta_j<0\).
The defining cusp average gives \(a_0(0)=0\).

The smooth periodic function at \(y=0\) has absolutely summable
Fourier coefficients: repeated integration by parts with
\(1-\sum_j\partial_{t_j}^2\) bounds them by
\(C_N(1+|\eta|^2)^{-N}\) for every \(N\).
At each \(y\ge0\) the resulting series represents \(A(t+iy)\)
by Fourier uniqueness. On the diagonal ray its nonzero modes
therefore tend uniformly to zero as the height tends to infinity,
by dominated convergence since \(\sum_j\eta_j>0\) on their support.
Finally \(G(\zeta)=A((\zeta/(k_0c_1^2))_\tau)\), since the
rank-one operator of \(c_1v\) is \(c_1^2\) times that of \(v\).
This proves the required uniform vanishing and the claimed
one-variable constant-term identity.

This exponential identifies the half-plane above
the boundary modulo
periods with a punctured disk up to radius \(e^{2\pi t/H'}\).
The induced function extends with value zero at \(q=0\). Apply
the maximum principle to its quotient by \(q\) and evaluate at
\(q=1\). This would give
\[
  1\le 2^{1+w/2}\exp(-2\pi t/H')
       \le 2^{1+w/2-(fw+2)}<1 .
\]
Thus the cusp test implies the required minimum. Its application
on the actual deep covers needs the same period bound on every
component for the lattice used in the return count.

The period has a direct lattice calculation once the integral model
is identified. In the host normalization put
\(c=\mathrm N(\mathfrak c_K)>0\); the purely imaginary \(\xi_0\)
satisfies \(|\tau(\xi_0)|=1/(c\sqrt d)\), and the transvection
coefficient \(\kappa_0=\pm2\xi_0\) satisfies
\((cd\kappa_0)^2=-4d\). Thus \(cd\kappa_0\) is integral. Suppose
the principal level is imposed on the lattice
\(\Lambda=\mathfrak a\Lambda'\), as in the actual comparison, with
\(\Lambda'\) Hermitian self-dual. For \(v\in\Lambda'\) isotropic
and \(y\in\Lambda\), \(h(y,v)v\in\Lambda\), using the convention
that \(h\) is linear in its first argument. The transvection
\[
        y\longmapsto y+Mcd\kappa_0\,h(y,v)v
\]
therefore preserves \(\Lambda\) and is the identity modulo \(M\).
It is unitary by \(\bar\kappa_0=-\kappa_0\) and \(h(v,v)=0\);
its inverse is obtained by changing the sign. Its raw period is
\(Mcd\), and scaling by \(k_0=2/(c\sqrt d)\) gives precisely
\(H=2M\sqrt d\). The same rational transformation realizes the
period simultaneously at every \(\tau\), since the imaginary
coefficient has the same negative imaginary part \(-k_0\) there.
For the even model weights the required arithmetic
character on these transvections must be trivial. The identification
of these lattices with the actual congruence groups on every deep
component, and the identification of the sections under test with
the stated analytic cusp space, remain necessary for this application.
\src{src/Development/AbelianVarietyTheory/Proofs/CubicalLineBundles/IsIntegralLevelMulKappa.lean}{35}{Integral transvection in the principal congruence subgroup}.

In the total mass comparison, the more precise intermediate upper
bound has the form
\[
   V\le 16^{fr^2} M^{2fr^2}\operatorname{vol}(U(r))^f
                            \operatorname{mass}_I .
\]
Here the compact group volume uses the real normalization in the
adelic tube comparison and \(\operatorname{mass}_I\) is the
definite standard lattice mass for the host. A bound
\(\operatorname{vol}(U(r))\le13^{r^2}\) and
\(\operatorname{mass}_I\le[3(2f)^2r^4R_E^2]^{6fr^2}\)
give \(V\le V_*\) since \(16\cdot13\le15^2\).
The principal-level index factor \(M^{2fr^2}\) itself follows by
reducing integral lattice automorphisms modulo \(M\). Their image
lies in the endomorphism ring of \(\Lambda/M\Lambda\), of order
\(M^{2fr^2}\) by local freeness of rank \(r\) over \(O_E\).
It remains to compare the adelic component tubes, with their disjointness
and volumes, and to identify the pushforward probability in the second
hypothesis. We calculate the compact factor and the definite mass below.
\src{src/ReciprocityTheory/ReductiveGroups/ArithmeticGeometry/FundamentalDomains/MaxBoundPf.lean}{12}{Isotropic flow and periodic cusp restriction};
\src{src/Development/WeilBound/RepresentationTheoryGL/TubeMeasures/DisjointTubes/MeasureSumBounds/SumMuFLe.lean}{110}{Principal level tube comparison for all components};
\src{src/Development/AutomorphicArithmetic/ThetaCorrespondence/FockModels/TamagawaPeterssonLambdaSqNonneg.lean}{371}{Restriction to normalized algebraic components}.

\paragraph{The compact real factor and the definite lattice mass.}
Put \(n=[E:\mathbb Q]=2f\), and write \(O=O_E\). The two numerical
factors in the preceding comparison have the following bounds.
\begin{proposition}\label{bergvol:compactmass}
For the real density specified below and the genus of the standard
positive definite Hermitian \(O\)-lattice of rank \(r\ge1\),
\[
 c_r=\operatorname{vol}(U(r))
     =\prod_{j=1}^r\frac{2\pi^j}{(j-1)!}\le13^{r^2},
 \qquad
 \operatorname{mass}_I
     \le [3n^2r^2R_E^2]^{2nr^2}
     \le[3n^2r^4R_E^2]^{3nr^2}.
\]
\end{proposition}
\begin{proof}
In the skew-Hermitian Lie algebra use the density \(dx\) given by
the imaginary diagonal entries and the real and imaginary parts of
the entries strictly below the diagonal. To compare it with the real
measure of the tube calculation, let \(\lambda\) be the additive
measure for which \((X,Y)\mapsto X+iY\) carries
\(\lambda\otimes\lambda\) to entrywise Lebesgue measure on
\(M_r(\mathbb C)\). Each diagonal block has absolute real determinant
one. Each off-diagonal pair gives the two complex entries
\(z+iw,-\bar z-i\bar w\), a block of absolute real determinant four.
Hence \(\lambda=2^{\binom r2}dx\). The specified real Haar density
at the identity, in exponential coordinates, is
\(2^{-\binom r2}\lambda=dx\).
\src{src/Formalism/LocalRingCardinality/DataCD.lean}{42}{Complexification measure on the real Lie algebra};
\src{src/Formalism/LocalToGlobal/BesselHermitianTheta/MinfDensity.lean}{13}{Real density relative to complexification}.

The map \(U(j)\to S^{2j-1}\), \(g\mapsto ge_1\), has fibre
\(U(j-1)\). At the identity its differential selects the first
column, whose independent coordinates give Euclidean tangent
measure on the sphere at \(e_1\). The remaining coordinates give
exactly the chosen density on \(U(j-1)\). By invariance the same
quotient measure is Euclidean surface measure on the whole sphere.
Starting with \(c_0=1\) gives the displayed product. For \(j\le r\)
each factor is at most \(2\cdot4^r\le8^r\).
In particular compact probability Haar has density
\(2^{-\binom r2}/c_r\) relative to \(\lambda\). For a real form
\(U(p,q)\), \(p+q=r\), the compact-dual convention prescribes this
same scalar relative to its own complexification measure. Thus
multiplication by \(c_r\) converts that convention to the specified
real density at each place.
\src{src/Formalism/LocalRingCardinality/DataCD.lean}{57}{Matching of the real density across signatures};
\src{src/ReciprocityTheory/ReductiveGroups/DiscreteSeries/FormalDegree/VolUFactProdMulLe.lean}{14}{Compact real volume factor}.

For the second assertion, every lattice \(L\) in the standard genus
is Hermitian self-dual. Indeed this holds locally for \(O^r\), and
localization commutes with the dual of a projective lattice. One may
use the conjugation-stable semilocal rings at pairs
\(\{\mathfrak p,\bar{\mathfrak p}\}\) here, including at split primes.
We bound the number of isometry classes of all totally positive
definite self-dual Hermitian \(O\)-lattices of rank \(r\).
Give the underlying real space the scalar product
\(\tfrac12\operatorname{Tr}_{E/\mathbb Q}h(x,y)\). Its squared norm
is \(\sum_\tau h_\tau(x_\tau,x_\tau)\), with one summand per conjugate
pair. For the absolute different \(\mathfrak D\) the Euclidean
dual lattice is \(L^*=2\mathfrak D^{-1}L\). Indeed testing the
pairing on all \(O\)-multiples of each lattice vector gives precisely
the codifferent condition, and \(L\) is Hermitian self-dual.
For \(0\ne y\in\mathfrak D^{-1}L\) we have
\(h(y,y)\in\mathfrak D^{-1}\bar{\mathfrak D}^{-1}\), whence
\(\lvert N_{E/\mathbb Q}h(y,y)\rvert\ge\Delta^{-2}\).
The arithmetic-geometric mean inequality gives
\(\|y\|^2\ge f\Delta^{-2/n}=fR_E^{-2}\).
Thus \(\lambda_1(L^*)\ge 2\sqrt f/R_E\).
By Banaszczyk's Euclidean transference theorem in dimension \(nr\),
\(\lambda_{nr}(L)\lambda_1(L^*)\le nr\)
\cite[Theorem~2.1]{banaszczyk}. Among real-linearly independent
vectors realizing this last minimum we may choose an \(E\)-basis
\(x_1,\ldots,x_r\) with \(\|x_i\|\le T=nrR_E\).

Write \(G=(h(x_i,x_j))\) and \(L_0=\sum O x_i\). The entries of
\(G\) are integral and have modulus at most \(T^2\) at every
embedding. There are at most \((1+2T^2)^n\) choices per entry.
Indeed the points of \(O\) in its Minkowski space have mutual sup distance
at least one by the norm inequality, so open polydiscs of radius
\(1/2\) around these points are disjoint. For a fixed \(G\) the marked
rational Hermitian space and \(L_0\) are fixed up to isometry.
The inclusions \(L_0\subseteq L\subseteq L_0^\#\), where \(\#\)
denotes the Hermitian dual, locate \(L/L_0\) in a finite module
\(A=L_0^\#/L_0\) of order \(\lvert N_{E/\mathbb Q}\det G\rvert
\le T^{2nr}\). The last inequality follows by Gram's determinant
inequality at every embedding.

Each \(O\)-submodule of \(A\) has an ordered generating tuple of
length \(r\). Locally at each prime in the support its inverse image
in \((L_0^\#)_{\mathfrak p}\) is free of rank \(r\) over a discrete
valuation ring, so its image has \(r\) generators. Combine these
tuples coordinate by coordinate under the primary decomposition of
the finite module. Consequently there are at most \(|A|^r\)
intermediate lattices for the fixed Gram matrix. All isometry
classes together number at most
\[
 (1+2T^2)^{nr^2}T^{2nr^2}\le (3T^2)^{2nr^2}.
\]
The definite standard mass is the sum of
\(1/|\operatorname{Aut}_O(L,h)|\) over the classes in its genus.
Each group is finite, since it is a discrete lattice stabilizer
inside a compact real group. The class count therefore proves
the required bound, including for nonfree \(L\).
\src{src/Formalism/SharpRestriction/GlobalSchurBounds/GeneralLinearActions/GapAtLeast.lean}{20}{Definition of the definite standard mass};
\src{src/Formalism/HermitianLattices/StatementMass.lean}{11}{Reciprocal stabilizer weights}.
\end{proof}

These estimates give \(V\le V_*\) once the displayed adelic tube
inequality holds. Applying that inequality to a different signature
still requires the comparison of the adelic covolumes with their
finite lattice stabilizers and the stated real densities.


\begin{corollary}[Size of the component bound]\label{bergvol:size}
Suppose \(L\ge1,\ \log(fr)\le L,\ \log M\le30L,\ \log R_E\le3L,\)
\(d^f\le\Delta\), and \(\log w\le25L\). Then
\[
  1\le\beta_*,\qquad \log\beta_*\le 400 f^2 r^2 L
          \le4000(frL)^3 .
\]
\end{corollary}
\begin{proof}
We have \(\log d\le6L\), \(\log\Delta\le6fL\),
\(\rho^{-1}\le e^{35L}\), and \(\varepsilon^{-1}\le
3fwM\sqrt d\le e^{61L}\). The radius estimate uses
\(\log(w+f)\le\log2+25L\); the bound on \(\varepsilon^{-1}\)
uses \(fw+2\le3fw\) and \(2\log2\le\pi\).
Since \(1\le b\le fr^2/2\) and \(\log b\le2L\), the logarithm
of \(8b!/(\pi\rho^2)^b\) is at most \(3L+36fr^2L\).
The logarithm of the base in \(A\) is at most
\(\log16+2\log(fr)+\log\Delta+\log(\varepsilon^{-1})
\le(66+6f)L\).
In \(V_*\), the first factor contributes at most \(66fr^2L\);
the logarithm of the base of its second factor is bounded by
\(\log3+2\log(2f)+4\log r+2\log R_E\le16L\).
Adding gives \(\log\beta_*\le400 f^2r^2L\). Each factor in
\(\beta_*\) is at least one, in particular \(\pi\rho^2<1\).
The final inequality uses \(f,r,L\ge1\).
\end{proof}
For the cover weight take \(w=2[E:\mathbb Q]\sum_{\tau,i}
|\lambda_{\tau i}|+1\). If the weight sum is zero then \(w=1\).
Otherwise it is an integer at least one; if its logarithm is at
most \(20L\) and \(\log[E:\mathbb Q]\le3L\), then
\(\log w\le\log3+23L\le25L\). The stage's field size
\(\log[E:\mathbb Q]+\log R_E\le3L\) implies both field estimates.
\src{src/Development/AutomorphicArithmetic/ThetaCorrespondence/FockModels/TamagawaPeterssonLambdaSqNonneg.lean}{458}{Cubic cost with the host discriminant bound};
\src{src/Development/AutomorphicArithmetic/ThetaCorrespondence/ThetaPairLifts/HodgeBundleWeights/IntegralCuspSlopePowerBound.lean}{137}{Use of the same bound for the component kernel}.
